\documentclass[12pt]{article}
\usepackage[utf8]{inputenc}
\usepackage[T1]{fontenc}
\usepackage{hyperref}
\title{Mwaghavul Language Resources for K-12 Curriculum Development}
\author{Compiled by Claude Sonnet 5.5 (claude-sonnet-5-5) for Andr\'as Kornai}
\date{}
\begin{document}
\maketitle

\section{Introduction}
Mwaghavul (also Sura; ISO 639-3 \texttt{sur}) is a West Chadic language spoken mainly around Mangu and Panyam in Plateau State, Nigeria. Estimates of speaker numbers range from about 150,000 (Ethnologue 2016) to about 200,000 (Blench 2020, informal), below the 300,000 in the task description. Missionaries first wrote the language in the early 20th century using Yoruba-based spelling, and a new orthography was developed in the early 21st century.

\section{Primary Educational Resources}
No primers, textbooks or school programmes were found in search results. A Blench dictionary reportedly aims at children of inter-tribal marriages and those raised outside the chiefdom, but I did not find its URL in the results, so it is not listed as a resource.

\section{Linguistic and Reference Materials}
\begin{itemize}
\item Wikipedia, ``Mwaghavul language'': \url{https://en.wikipedia.org/wiki/Mwaghavul_language}
\item Omniglot page on Mwaghavul writing: \url{https://omniglot.com/writing/mwaghavul.htm}
\item Texas State University Pressbooks, ``Mwaghavul (Sura)'' chapter: \url{https://pressbooks.txst.edu/minoritylanguages/chapter/introduction-2/}
\item Mwaghavul minority-language project paper (Texas State, PDF): \url{https://pressbooks.txst.edu/app/uploads/sites/4/2023/02/Mwaghavul-Minority-Minoritized-Languages-and-Cultures-Project-AA.pdf}
\end{itemize}

\section{Religious and Community Resources}
\begin{itemize}
\item Joshua Project, Mwaghavul in Nigeria: \url{https://joshuaproject.net/people_groups/15128}
\end{itemize}

\section{Digital and Audio Resources}
\begin{itemize}
\item Global Recordings Network, Mwaghavul (audio recordings; the page notes a New Testament exists): \url{https://globalrecordings.net/en/language/sur}
\end{itemize}

\section{Curriculum Development Context}
The language is spoken widely, and a Bible tradition and an updated orthography exist. No school use was found. The main gap is classroom material, so developers would need to build on the new orthography and any dictionary work by R. Blench, which should be located and checked directly.

\end{document}